\documentclass[12pt]{beamer}
\usetheme{Pittsburgh}
\usecolortheme{seagull}
\usepackage[utf8]{inputenc}
\usepackage[english]{babel}
\usepackage{amsmath}
\usepackage{listings}
\usepackage{amsfonts}
\usepackage{amssymb}
\usepackage{graphicx}
\author{Introdução ao C++}
\title{Estruturas de Dados}
%\setbeamercovered{transparent} 
\setbeamertemplate{navigation symbols}{} 
%\logo{\includegraphics[scale=0.015]{Brasao_UFSC.png}\includegraphics[scale=0.2]{brasao_PPGCC.jpg}} 
\institute{Departamento de Informática e de Estatística \\ Prof. Jean Everson Martina \\ Prof. Aldo von Wangenheim} 
\date{2016.2} 
\subject{} 
\usebackgroundtemplate{\includegraphics[width=\paperwidth,
height=\paperheight]{../reusable_images/fundo_UFSC.png}}
\begin{document}

{
\usebackgroundtemplate{\includegraphics[width=\paperwidth,
height=\paperheight]{../reusable_images/fundo_capa.png}}
\begin{frame}
\titlepage
\includegraphics[scale=0.3]{../reusable_images/brasao_INE.png}
\end{frame}
}

\begin{frame}{Visão Geral}
\begin{itemize}
\item Histórico da Linguagem;
\item Programa C++: header, source – função main();
\item GCC/G++;
\item Arquivos objeto, bibliotecas dinâmicas e estáticas;
\item \#include, \#define, namespaces, typedefs;
\item Ponteiros, Referências, arrays, primitivas, estruturas de controle;
\end{itemize}
\end{frame}

\begin{frame}{Características do C++}
\begin{itemize}
\item Linguagem multi-paradigma;
\item Programação Procedural;
\item Orientação a objetos;
\item Meta Programação;
\item Desalocação memória – Decisão cabe ao programador;
\end{itemize}
\end{frame}


\begin{frame}{Histórico}
\begin{itemize}
\item Criada por Bjarne Stroustrup em 1979 no Bell Labs da AT\&T e inicialmente chamada de ``C with classes”;
\item Em 1983 passou a ser chamada C++;
\item Ao longo do tempo incorporou novas funcionalidades como:
\begin{itemize}
\item Herança múltipla; 
\item Sobreescrita de operadores;
\item Templates;
\item Entre outras coisas que veremos ao longo do curso.
\end{itemize}
\end{itemize}
\end{frame}

\begin{frame}{Histórico}
\begin{itemize}
\item Possui compatibilidade com C;
\item Um compilador C++ deve compilar um programa escrito em C;
\item Existem vários comiladores da linguagem;
\item É um linguagem padronizada internacionamelmente:
\begin{itemize}
\item Em 1998 foi definido um padrão ISO para a linguagem; 
\item Revisado em 2003;
\item Revisado em 2011;
\item Revisado em 2014;
\end{itemize}
\end{itemize}
\end{frame}

\begin{frame}{GCC/G++}
\begin{itemize}
\item GNU Compiler Collection;
\item Compilador originalmente desenvolvido focando a linguagem C, chamava-se GNU C 
Compiler;
\item Ao longo do anos passou a contemplar outras linguagens, como C++;
\end{itemize}
\end{frame}

\begin{frame}{Programa C++ - Arquivos}
\begin{itemize}
\item Header - Definições visíveis a outros headers e ao compilador;
\item Sources - Corpo de métodos e funções;
\item Implemented Headers - Headers com implementado que tem que ser instrumentada pelo compilador;
\item Função Main - Função inicial invocada pelo executável;
\end{itemize}
\end{frame}

\begin{frame}[fragile]{Programa exemplo}
\lstset{language=C++,
          keywordstyle=\color{blue}\ttfamily,
          stringstyle=\color{red}\ttfamily,
          commentstyle=\color{green}\ttfamily,
          breaklines=true,
          basicstyle=\ttfamily\footnotesize
          }
\begin{lstlisting}
#include <iostream>  // std::cout
int main()
{
    // Um programa bastante simples em C++, utilizando
    // tipos integrais e texto, e imprimindo 
    // algumas informacoes na tela:
    int year = 2016;
    unsigned age = 23;
    std::string name = "Tarcisio";

    std::cout << "Hello, " << name
              << ". I see you were born in "
              << year - age << "!\n";
    return 0;
}
\end{lstlisting}
\footnote{https://gitlab.com/tarcisioe/cpp/}
\end{frame}

\begin{frame}{Programa exemplo}
\begin{itemize}
\item Como podemos ver, começamos o programa com duas variáveis;
\item Variáveis em C++ possuem um tipo e um valor;
\item O tipo de uma variável nunca muda, mas seu valor pode mudar durante a execução do programa;
\item A variável `year` é uma variável do tipo `int`, que é utilizado para representar números inteiros com sinal;
\end{itemize}
\end{frame}

\begin{frame}{Programa exemplo}
\begin{itemize}
\item A variável `age`, por sua vez, representa uma idade, que não pode ser negativa. Por isso, utilizamos o tipo `unsigned`; 
\item A variável `name`, por sua vez, é de um tipo não-primitivo, ou seja, um tipo definido em C++, em algum lugar da biblioteca padrão;
\item Após as variáveis, a linha iniciada em `std::cout` utiliza o operador `<<` para escrever coisas na tela com a representação que C++ utiliza.
\end{itemize}
\end{frame}

\begin{frame}{Realização dos Trabalhos}
\begin{itemize}
\item Os trabalhos pequenos são todos individuais;
\item Equipes de 2 pessoas (somente trabalhos grandes):
\begin{itemize}
\item É importante aprender a realizar trabalhos de implementação em equipe;
\end{itemize}
\item Implementação:
\begin{itemize}
\item ANSI C++ - OOP;
\item Plataforma homologada é somente o VPL;
\end{itemize}
\item Ferramentas \& Máquinas:
\begin{itemize}
\item Aproveite as Aulas de Laboratório;
\item Linux e Eclipse para C/C++;
\item Compiladores: gcc / g++
\end{itemize}
\end{itemize}
\end{frame}

\begin{frame}{Entrega dos Trabalhos}
\begin{itemize}
\item Data:
\begin{itemize}
\item Deverão ser entregues na data especificada no Moodle;
\item São aceitos com atraso mas com 50\% de desconto na nota por atraso;
\item Projetos de implementação também;
\end{itemize}
\item Serão aceitos os trabalhos que:
\begin{itemize}
\item Satisfizerem as especificações / enunciados dados;
\item Utilizarem a linguagem e a técnica de programação especificadas;
\item Utilizarem codificação de acordo com os padrões internacionais para a linguagem especificada;
\item Estiverem devidamente documentados.
\item Que executem na plataforma homologada.
\end{itemize}
\end{itemize}
\end{frame}

\begin{frame}{Programa da Disciplina}
\begin{itemize}
\item Data:
\begin{itemize}
\item Deverão ser entregues na data especificada no Moodle;
\item São aceitos com atraso mas com 50\% de desconto na nota por atraso;
\item Projetos de implementação também;
\end{itemize}
\item Serão aceitos os trabalhos que:
\begin{itemize}
\item Satisfizerem as especificações / enunciados dados;
\item Utilizarem a linguagem e a técnica de programação especificadas;
\item Utilizarem codificação de acordo com os padrões internacionais para a linguagem especificada;
\item Estiverem devidamente documentados;
\item Que executem na plataforma homologada.
\end{itemize}
\end{itemize}
\end{frame}

\begin{frame}{Controle de Freqüência}
\begin{itemize}
\item O controle de freqüência segue as regulamentações da UFSC para cursos presenciais;
\item É obrigatório 75\% de presença para aprovação;
\item O controle de freqüência vai ser feito diretamente no moodle da disciplina;
\item O controle de freqüência poderá ser feito de forma automatizada usando a carteirinha do RU;
\begin{itemize}
\item Se você ainda não tem uma é so ir na secretaria do RU e emitir a sua.
\end{itemize}
\end{itemize}
\end{frame}

\begin{frame}{Programa da Disciplina}
\begin{itemize}
\item Introdução e Aspectos Gerais da Disciplina;
\item Introdução à Programação C++;
\item Listas, Pilhas e Filas utilizando Vetores (arrays);
\item Gerência e Alocação Dinâmica de Memória;
\item Lista Encadeada;
\item Fila e Pilha Baseados em Listas Encadeadas;
\item Lista Duplamente Encadeada;
\item Lista Circular;
\end{itemize}
\end{frame}

\begin{frame}{Programa da Disciplina}
\begin{itemize}
\item Conceitos de Complexidade de Algoritmos;
\item Introdução a Árvores;
\item Árvores I;
\item Árvores de Busca II;
\item Hashing;
\item Gerência de Arquivos;
\item Ordenação;
\end{itemize}
\end{frame}

\begin{frame}{Bibliografia}
\begin{enumerate}
\item Material de Estruturas de Dados:
\begin{itemize}
\item Horowitz: Estruturas de Dados. Ed. Campus;
\item Wirth: Data Structures + Algorithms = Programs;
\item Outros livros, a maioria em inglês, todos disponíveis na BU ou no INE;
\end{itemize}
\item Material sobre  C++:
\begin{itemize}
\item Livros sobre Estruturas de Dados em “C” e C++;
\item Material sobre linguagem C++;
\end{itemize}
\item Bibliografia de Ferramentas de Programação:
\begin{itemize}
\item Livros sobre Unix / Linux;
\item Manuais de Editores;
\item Manuais de Programação;
\item Material Online.
\end{itemize}
\end{enumerate}
\end{frame}

\begin{frame}{Contato e Marcação de Horários}
\begin{itemize}
\item Monitoria: será disponibilizado um horário por semana para atendimento pelo estagiário de docência ou bolsista da disciplina;
\item Procedimento: 
\begin{itemize}
\item Envie e-mail explicando seu problema;
\item Utilize sempre o Fórum de Discussão;
\item Se possível, o problema será resolvido por e-mail;
\item Caso contrário, será indicado o horário de atendimento.
\end{itemize}
\item Marcação de Horários
\begin{itemize}
\item Professores por mensagem no moodle;
\item Estagiário de docência e bolsista por e-mail divulgado no moodle;
\end{itemize}
\end{itemize}
\end{frame}

{
\usebackgroundtemplate{\includegraphics[width=\paperwidth,
height=\paperheight]{../reusable_images/fundo_capa.png}}
\begin{frame}

{\LARGE Perguntas????}

\end{frame}
}


{
\usebackgroundtemplate{\includegraphics[width=\paperwidth,
height=\paperheight]{../reusable_images/fundo_capa.png}}
\begin{frame}
\includegraphics[scale=0.8]{../reusable_images/cc_logo_arge.png}\hspace{0.5cm} 
\includegraphics[scale=0.95]{../reusable_images/by.png}

\vspace{1cm}
Este trabalho está licenciado sob uma Licença Creative Commons Atribuição 4.0 Internacional. Para ver uma cópia desta licença, visite http://creativecommons.org/licenses/by/4.0/.

\end{frame}
}
\end{document}